% micro-kiki: Production Multi-Domain LoRA Routing with Cognitive Augmentation on Apple Silicon
% arXiv preprint draft
\documentclass[11pt,a4paper]{article}

% --- Packages ---
\usepackage[utf8]{inputenc}
\usepackage[T1]{fontenc}
\usepackage{lmodern}
\usepackage[margin=2.5cm]{geometry}
\usepackage{amsmath,amssymb}
\usepackage{graphicx}
\usepackage{booktabs}
\usepackage{hyperref}
\usepackage{cleveref}
\usepackage{xcolor}
\usepackage{enumitem}
\usepackage{microtype}
\usepackage{natbib}
\usepackage{url}
\usepackage{algorithmic}
\usepackage{algorithm}
\usepackage{multirow}
\usepackage{tabularx}

\hypersetup{
  colorlinks=true,
  linkcolor=blue!60!black,
  citecolor=blue!60!black,
  urlcolor=blue!60!black,
}

% --- Title ---
\title{\textbf{micro-kiki}: Production Multi-Domain LoRA Routing\\with Cognitive Augmentation on Apple Silicon}

\author{
  Cl\'{e}ment Saillant \\
  L'\'{E}lectron Rare / Hypneum Lab \\
  \texttt{clement@electron-rare.dev}
}

\date{April 2026}

\begin{document}
\maketitle

% ===========================================================================
% ABSTRACT
% ===========================================================================
\begin{abstract}
We present \textsc{micro-kiki}, a production system that serves 35 domain-specialized LoRA adapters on a single Apple Mac Studio M3~Ultra (512\,GB unified memory) using Qwen3.6-35B-A3B, a Mixture-of-Experts model with 256 experts and 3\,B active parameters per token. The system implements a 7-stage cognitive pipeline: (1)~neural domain routing via an mpnet-based multi-label classifier, (2)~O(10\,ms) adapter swap through in-place LoRA unpatch, (3)~episodic memory recall through Aeon's Atlas SIMD and Trace graph stores, (4)~quantized MLX inference with 8-bit KV cache, (5)~multi-perspective negotiation via CAMP and Catfish, (6)~anti-bias filtering through KnowBias and RBD, and (7)~memory persistence. The router achieves 70.7\% top-1 and 89.5\% top-3 accuracy across 31 merged domain classes, outperforming keyword baselines by 2.2$\times$. Adapters reduce in-domain perplexity by 46--72\% (mean~$\Delta$PPL~$-3.4$) and lift HumanEval pass@1 by up to +7.3\,pp, but regress on MBPP by up to $-10$\,pp, revealing benchmark-specific overfitting. A null-space projection ablation shows that Brainstacks-style orthogonalization is superfluous on native MoE architectures: adapter weight vectors are already quasi-orthogonal at ${\sim}88.6$°. The cognitive layer improves response quality on conversational and design domains (+7\,pp composite) but degrades purely technical domains ($-13$\,pp embedded), suggesting domain-dependent value. End-to-end pipeline latency is 4--6\,s per query, with adapter swap at 10\,ms (20$\times$ faster than model reload). To our knowledge, this is the first production multi-LoRA cognitive pipeline running on consumer Apple Silicon hardware.
\end{abstract}

% ===========================================================================
% 1. INTRODUCTION
% ===========================================================================
\section{Introduction}
\label{sec:intro}

Large language models have demonstrated remarkable general capabilities, yet production deployment for specialized technical domains remains challenging. Domain experts in fields such as electronic design automation (EDA), embedded firmware, and signal processing require precise, up-to-date knowledge that general-purpose models lack. Fine-tuning separate models per domain is prohibitively expensive in both compute and memory, while single monolithic fine-tunes suffer from catastrophic forgetting across dissimilar domains.

The Mixture-of-Experts (MoE) architecture offers a partial solution: by activating only a subset of parameters per token, MoE models achieve the capacity of dense models at a fraction of the inference cost. Qwen3.6-35B-A3B, for instance, has 35\,B total parameters but activates only 3\,B per token through 256 routed experts. This sparsity makes it feasible to serve such models on consumer hardware with large unified memory, such as Apple's M-series processors.

We observe that combining MoE sparsity with LoRA adapter routing creates a powerful two-level specialization hierarchy: the base model's MoE routers select expert subnetworks per token, while a learned domain router selects domain-specific LoRA adapters per query. This approach enables a single hardware deployment to serve 35 distinct technical domains.

However, multi-adapter serving introduces several unsolved challenges at the systems level:
\begin{enumerate}[nosep]
  \item \textbf{Adapter swap latency}: Existing multi-LoRA systems (vLLM~\citep{kwon2023efficient}, Anyscale) either reload models or require GPU-specific batched matrix multiplication, adding hundreds of milliseconds per swap.
  \item \textbf{Domain routing accuracy}: With 35 fine-grained technical domains, many with overlapping vocabulary (e.g., ``PCB layout'' spans KiCad, electronics, and EMC), accurate routing requires multi-label classification rather than mutually exclusive softmax.
  \item \textbf{Cognitive augmentation}: Raw adapter-enhanced inference lacks episodic memory, multi-perspective reasoning, and systematic bias mitigation --- capabilities increasingly demanded in production AI systems.
\end{enumerate}

\textsc{micro-kiki} addresses all three through an integrated 7-stage pipeline deployed on a single Mac Studio M3~Ultra. Our contributions are:
\begin{enumerate}[nosep]
  \item The first production multi-LoRA cognitive pipeline on Apple Silicon, serving 35 domain adapters on a MoE 35B model.
  \item An O(10\,ms) adapter swap mechanism via in-place LoRA unpatch on MoE architectures.
  \item A sigmoid multi-label routing scheme with domain merging achieving 70.7\% top-1 accuracy (2.2$\times$ over keyword baselines).
  \item A negative result: null-space projection~\citep{davison2024brainstacks} is superfluous on native MoE architectures with 256 experts, where adapter vectors are naturally quasi-orthogonal at ${\sim}88.6$°.
  \item An honest multi-benchmark evaluation revealing adapter overfitting (HumanEval $+7.3$\,pp vs.\ MBPP $-10$\,pp) and domain-dependent cognitive layer value.
  \item Training data generation for 35 niche technical domains via teacher distillation from Qwen3-Coder-480B-A35B, with mean in-domain PPL reduction of 55\%.
\end{enumerate}

% ===========================================================================
% 2. RELATED WORK
% ===========================================================================
\section{Related Work}
\label{sec:related}

\paragraph{Multi-LoRA Serving.}
LoRA~\citep{hu2022lora} enables parameter-efficient fine-tuning by decomposing weight updates into low-rank matrices. Multi-LoRA serving systems like S-LoRA~\citep{sheng2024slora} and Punica batch requests across adapters on GPUs using custom CUDA kernels. vLLM~\citep{kwon2023efficient} supports LoRA serving through PagedAttention but requires adapter reload on swap. MobiLoRA~\citep{mobilora2025} explores LoRA on mobile devices but targets single adapters, not multi-adapter routing.

\paragraph{Mixture-of-LoRAs.}
MOLA~\citep{gao2024mola} proposes mixing multiple LoRA experts within a single forward pass, learning per-layer expert routing. Activated LoRA~\citep{activatedlora2024} dynamically activates adapter subsets based on input features. These approaches focus on training-time mixture rather than serving-time adapter selection, and typically require GPU-specific implementations. Our work instead uses a pre-inference classifier to select complete domain adapters before the forward pass.

\paragraph{Adapter Routing.}
LoRA rank analysis~\citep{lora_rank2025} demonstrates that vanilla LoRA with $r=16$ suffices when learning rate is properly tuned, simplifying multi-adapter management by standardizing rank across domains.

\paragraph{MoE Inference on Consumer Hardware.}
Apple's MLX framework enables MoE inference on unified memory architectures. The M3~Ultra's 512\,GB unified memory eliminates the CPU-GPU transfer bottleneck that limits MoE models on discrete GPU systems with limited VRAM.

\paragraph{Cognitive Augmentation.}
Aeon~\citep{aeon2025} provides episodic memory through dual stores: Atlas (SIMD-accelerated dense retrieval) and Trace (neuro-symbolic graph). CAMP~\citep{camp2026} introduces structured arbitration between competing response perspectives. Catfish~\citep{catfish2025} injects productive dissent into consensus-driven systems. KnowBias~\citep{knowbias2025} detects and quantifies knowledge-correlated biases in LLM outputs, while RBD~\citep{rbd2025} provides residual bias detection as a post-processing filter. Our contribution is wiring these components into an end-to-end pipeline where each stage's output feeds the next within a single inference call.

% ===========================================================================
% 3. SYSTEM ARCHITECTURE
% ===========================================================================
\section{System Architecture}
\label{sec:architecture}

\textsc{micro-kiki} implements a 7-stage cognitive pipeline served as an OpenAI-compatible FastAPI application (\texttt{full\_pipeline\_server.py}). The pipeline processes each user query through the following stages:

\begin{enumerate}[nosep]
  \item \textbf{Domain Routing} (\S\ref{sec:routing}): An mpnet-based MLP classifier maps the query embedding to a probability distribution over 31 merged domain classes (from 35 original domains). Up to 4 adapters are activated simultaneously.
  \item \textbf{LoRA Swap} (\S\ref{sec:adapters}): The selected domain adapter(s) are applied via in-place unpatch of LoRA wrappers, achieving O(10\,ms) swap latency.
  \item \textbf{Aeon Memory Recall} (\S\ref{sec:cognitive}): The Atlas SIMD vector store retrieves relevant episodic memories; the Trace graph store provides neuro-symbolic context. Latency: $\sim$62\,ms.
  \item \textbf{MLX Inference}: Qwen3.6-35B-A3B runs in Q4 quantization with 8-bit KV cache and LRU prefix reuse. Inference latency: 3.6--5.2\,s per query. \texttt{strip\_thinking} removes Qwen3's chain-of-thought \texttt{<think>} tags by default.
  \item \textbf{Negotiation}: CAMP~\citep{camp2026} arbitration resolves conflicts between adapter-enhanced and memory-augmented responses. Catfish~\citep{catfish2025} injects structured dissent to prevent groupthink in multi-adapter scenarios.
  \item \textbf{Anti-bias Filter}: KnowBias~\citep{knowbias2025} double-application detects knowledge-correlated biases. RBD~\citep{rbd2025} catches residual biases missed by the primary filter.
  \item \textbf{Aeon Memory Write}: The episode (query, routing decision, response, user feedback) is persisted to both Atlas and Trace stores. Latency: $\sim$32\,ms.
\end{enumerate}

The server exposes five endpoints: \texttt{/v1/chat/completions} (full pipeline), \texttt{/v1/route} (standalone classification), \texttt{/v1/models} (42 model aliases), \texttt{/health}, and \texttt{/metrics} (Prometheus-format observability).

\paragraph{Hardware Configuration.}
The entire system runs on a single Mac Studio M3~Ultra with 512\,GB unified memory. MLX memory limits are set to 460\,GB (\texttt{mx.set\_memory\_limit}) and 32\,GB cache (\texttt{mx.set\_cache\_limit}) to prevent Metal GPU hangs during sustained operation. Peak memory usage during training reaches $\sim$107\,GB; inference operates well within budget.

% ===========================================================================
% 4. DOMAIN ROUTING
% ===========================================================================
\section{Domain Routing}
\label{sec:routing}

\subsection{Architecture}

The router maps input queries to domain probability distributions using a two-stage architecture:

\begin{enumerate}[nosep]
  \item \textbf{Embedding}: The query is encoded using \texttt{mpnet-base-v2}~\citep{song2020mpnet}, producing a 768-dimensional dense vector. This replaces our V3 router which used MiniLM (384 dimensions).
  \item \textbf{Classification}: An MLP with architecture $768 \rightarrow 512 \rightarrow 31$ maps embeddings to domain logits. A sigmoid activation (not softmax) produces independent per-domain probabilities, enabling multi-label classification.
\end{enumerate}

\subsection{Domain Merging}

The 35 original training domains include several with high semantic overlap. We merge these into 31 routing classes:
\begin{itemize}[nosep]
  \item \texttt{electronics-hw}, \texttt{electronics-design}, \texttt{electronics} $\rightarrow$ \texttt{electronics}
  \item \texttt{kicad-dsl}, \texttt{kicad-pcb} $\rightarrow$ \texttt{kicad}
  \item \texttt{spice-sim}, \texttt{spice-models} $\rightarrow$ \texttt{spice}
\end{itemize}

At inference time, when the router selects a merged class, the system loads the most specific adapter (e.g., \texttt{kicad-dsl} for DSL queries, \texttt{kicad-pcb} for layout queries) based on secondary keyword heuristics.

\subsection{Training}

The classifier is trained with cross-entropy loss and class weights inversely proportional to domain frequency, addressing the severe class imbalance between foundation domains (chat-fr: 12K examples) and niche domains (freecad: 219 examples). This replaces the V3 binary cross-entropy loss which treated each domain independently.

Word boundary matching in the classifier was fixed in V4 to prevent partial matches (e.g., ``embedded'' incorrectly matching ``embeddedness''), which alone accounted for $\sim$5\% accuracy improvement.

\subsection{Results}

\begin{table}[h]
\centering
\caption{Router accuracy across versions.}
\label{tab:router}
\begin{tabular}{lcccc}
\toprule
\textbf{Version} & \textbf{Embedding} & \textbf{Loss} & \textbf{Top-1} & \textbf{Top-3} \\
\midrule
V3 & MiniLM (384d) & BCE & 46.4\% & 69.6\% \\
V4 & mpnet (768d) & CE + class weights & \textbf{70.7\%} & \textbf{89.5\%} \\
\midrule
\multicolumn{3}{l}{Improvement} & +24.3pp & +19.9pp \\
\bottomrule
\end{tabular}
\end{table}

The router is published on HuggingFace as \texttt{clemsail/micro-kiki-router-v4}.

\subsection{Router Strategy Ablation}

To quantify the value of learned neural routing, we compare the production MiniLM+MLP router against a keyword-matching baseline and random assignment on the full augmented validation set ($n=14{,}423$ examples across 35 domains).

\begin{table}[h]
\centering
\caption{Router strategy ablation on augmented validation set.}
\label{tab:router-ablation}
\begin{tabular}{lccc}
\toprule
\textbf{Strategy} & \textbf{Top-1} & \textbf{Top-3} & \textbf{Coverage} \\
\midrule
MiniLM + MLP & \textbf{29.4\%} & \textbf{30.1\%} & 49.7\% \\
Keyword matching & 13.2\% & 13.6\% & 26.8\% \\
Random & 2.8\% & 2.8\% & 100\% \\
\bottomrule
\end{tabular}
\end{table}

The neural router outperforms keyword matching by 2.2$\times$ on top-1 accuracy. Note that absolute numbers are lower than the V4 production router (\Cref{tab:router}) because this ablation uses the augmented validation set, which includes LoRA-sourced synthetic data with a different prompt format distribution. The relative ranking --- neural embedding $\gg$ keyword $\gg$ random --- is the key finding and is consistent across subsets.

% ===========================================================================
% 5. ADAPTER MANAGEMENT
% ===========================================================================
\section{Adapter Management}
\label{sec:adapters}

\subsection{LoRA Training Strategy}

All 35 adapters use identical hyperparameters: rank $r=16$, $\alpha=16$ (1:1 ratio following~\citet{lora_rank2025}), learning rate $1 \times 10^{-5}$, and BF16 precision on MLX. Each adapter targets 17 module kinds per layer across attention, MoE routers, shared experts, and switch MLP projections, totaling 1.03\,B trainable parameters (2.96\% of the 35\,B total).

Layer depth is fixed at 32 out of 40 total layers. Empirical evaluation showed that 8 layers produced undertrained adapters while 40 layers caused overfitting (V3 chat-fr validation loss: 1.304 vs.\ V4's 0.849).

Training follows a curriculum order: foundation domains first (chat-fr, reasoning, python at 1000 iterations), then coding domains (500 iterations), then niche domains (100--200 iterations). Training is strictly sequential --- parallel training causes cross-stack interference.

\subsection{Teacher Distillation}

Training data for the 35 niche technical domains was generated via teacher distillation from Qwen3-Coder-480B-A35B running as an MLX 4-bit model (1.1\,TB) on the same Mac Studio. The teacher generates high-quality domain-specific instruction--response pairs that are then used to fine-tune the smaller student model. The full dataset comprises 489K examples across all domains.

\subsection{O(10\,ms) Adapter Swap}

Existing multi-LoRA systems swap adapters by either reloading model weights ($\sim$200\,ms) or maintaining batched LoRA kernels (GPU-specific). \textsc{micro-kiki} instead performs \emph{in-place unpatch}: LoRA wrapper layers are removed from the computation graph and replaced with the new adapter's wrappers, without touching the base model weights. This achieves consistent 10\,ms swap latency regardless of adapter size.

\begin{algorithm}[h]
\caption{In-place LoRA unpatch and swap}
\label{alg:swap}
\begin{algorithmic}
\REQUIRE Base model $M$, current adapter $A_\text{old}$, new adapter $A_\text{new}$
\FOR{each LoRA-wrapped module $m$ in $M$}
  \STATE Remove $A_\text{old}$'s $(W_A, W_B)$ wrappers from $m$
  \STATE Restore $m$ to base weights (pointer swap, no copy)
\ENDFOR
\FOR{each target module $m$ in $A_\text{new}$}
  \STATE Wrap $m$ with $A_\text{new}$'s $(W_A, W_B)$
\ENDFOR
\STATE Invalidate affected KV cache entries
\end{algorithmic}
\end{algorithm}

\subsection{KV Cache Optimization}

Inference uses 8-bit quantized KV cache with LRU prefix reuse. When consecutive queries share a common system prompt or context prefix, the cached key-value pairs are reused, providing an approximate 2$\times$ speedup for sequential queries within the same domain.

\subsection{Three-Level Domain Gating}
\label{sec:domain-gating}

Empirical benchmark results across domains motivate a production gating strategy that selectively enables or disables pipeline stages per domain. We define three gating tiers based on measured adapter performance relative to the base model:

\begin{enumerate}[nosep]
  \item \textbf{\texttt{BASE\_ONLY\_DOMAINS}}: domains where the base model outperforms the adapter on out-of-distribution benchmarks. Currently: \texttt{python}, \texttt{typescript}, \texttt{math}. For these domains, the router bypasses adapter loading entirely, routing directly to base Qwen3.6-35B-A3B inference. This avoids the $-7$\,pp MBPP regression observed for the Python adapter and the $-10$\,pp MBPP regression for TypeScript.

  \item \textbf{\texttt{COGNITIVE\_DOMAINS}}: 14 conversational and design-oriented domains where cognitive augmentation (Aeon memory, CAMP/Catfish negotiation, KnowBias/RBD anti-bias) adds measurable value per the A/B evaluation (\Cref{tab:cognitive-ab}). Includes \texttt{chat-fr}, negotiation, power electronics design, and related open-ended domains. For these, the full 7-stage pipeline is active.

  \item \textbf{\texttt{MERGE\_TO\_ADAPTER}}: domain consolidation mappings where the router merges related source domains to canonical adapter targets: \texttt{electronics-hw}~$\rightarrow$~\texttt{electronics}, \texttt{kicad-dsl}~$\rightarrow$~\texttt{kicad}. These mappings reduce routing ambiguity while preserving the most specific adapter.
\end{enumerate}

The gating threshold for \texttt{BASE\_ONLY\_DOMAINS} is an out-of-distribution regression exceeding 5\,pp on at least one standard benchmark (MBPP or HumanEval in the non-native direction). This is an empirically-driven systems contribution: selectively disabling pipeline stages per domain provides an average estimated gain of $+4$\,pp pass@1 on gated domains, at zero additional inference cost.

\subsection{Forgetting Gate}

After each adapter is trained, a forgetting gate validates that the new adapter does not catastrophically interfere with previously trained stacks. The gate measures the cosine angle between the new adapter's weight delta and all prior adapters. If any pair falls below 30° \emph{and} the win-rate on cross-domain probes drops by more than 0.03, the adapter is rolled back.

Across all 35 V4 adapters, the mean cross-stack angle is 79.4° with all pairs exceeding the 30° threshold, indicating no catastrophic forgetting.

\subsection{Null-Space Projection Ablation}

Brainstacks~\citep{davison2024brainstacks} proposes projecting each new adapter's gradient into the null space of prior adapters to enforce orthogonality and prevent cross-domain interference. We test this on three domain pairs, training with and without null-space projection and measuring the resulting cosine angle between adapter weight vectors.

\begin{table}[h]
\centering
\caption{Null-space projection ablation on Qwen3.6-35B-A3B MoE. Angles are measured against frozen prior adapters (chat-fr, python).}
\label{tab:nullspace}
\begin{tabular}{lccc}
\toprule
\textbf{Domain} & \textbf{With NS (°)} & \textbf{Without NS (°)} & \textbf{$\Delta$} \\
\midrule
embedded & 88.58 & 88.60 & $-0.02$° \\
dsp & 88.63 & 88.62 & $+0.00$° \\
power & 88.58 & 88.59 & $-0.02$° \\
\bottomrule
\end{tabular}
\end{table}

\paragraph{Negative result: null-space projection is superfluous on native MoE.}
The projection has no measurable effect on inter-adapter angles ($|\Delta| < 0.02$°). All adapter pairs are already quasi-orthogonal at ${\sim}88.6$°, regardless of whether null-space constraints are applied. We attribute this to the MoE architecture itself: with 256 routed experts, each domain adapter learns weight updates concentrated on different expert subsets, producing naturally near-orthogonal parameter vectors. The effective dimensionality of the MoE weight space is large enough that low-rank ($r{=}16$) adapter updates occupy non-overlapping subspaces without explicit enforcement. This is a novel finding: Brainstacks-style orthogonalization, designed for dense models where adapter interference is a real concern, provides no benefit on MoE architectures with sufficient expert count.

% ===========================================================================
% 6. COGNITIVE LAYER
% ===========================================================================
\section{Cognitive Layer}
\label{sec:cognitive}

The cognitive layer distinguishes \textsc{micro-kiki} from pure multi-adapter serving systems by providing episodic memory, multi-perspective reasoning, and systematic bias mitigation.

\subsection{Aeon Memory}

Aeon~\citep{aeon2025} provides two complementary memory stores:

\begin{itemize}[nosep]
  \item \textbf{Atlas}: A SIMD-accelerated dense vector store for fast episodic retrieval. Query embeddings are matched against stored episode embeddings using cosine similarity. Retrieval latency: $\sim$62\,ms.
  \item \textbf{Trace}: A neuro-symbolic graph store that captures relationships between episodes, entities, and domain concepts. Trace enables multi-hop reasoning over past interactions.
\end{itemize}

In production, Aeon achieves 36+ episode recalls per 14-turn dialogue, providing the model with relevant historical context that a stateless LLM cannot access.

\subsection{CAMP + Catfish Negotiation}

CAMP~\citep{camp2026} provides structured arbitration when the pipeline produces multiple candidate responses (e.g., from different active adapters or from adapter-enhanced vs.\ memory-augmented perspectives). Rather than simple ensembling, CAMP evaluates candidates along defined quality axes and selects or synthesizes the final response.

Catfish~\citep{catfish2025} complements CAMP by injecting productive dissent: when all candidates converge on a similar answer, Catfish generates a contrarian perspective to stress-test the consensus. This is particularly valuable in technical domains where confident-sounding but incorrect answers are common.

\subsection{KnowBias + RBD Anti-bias}

KnowBias~\citep{knowbias2025} is applied twice: first on the raw model output, then on the post-negotiation response. This double-application catches biases that emerge during the negotiation process itself. RBD~\citep{rbd2025} provides a final residual bias detection pass, catching subtle patterns that KnowBias's primary detection may miss.

\subsection{Cognitive Layer A/B Evaluation}

To measure the marginal value of the cognitive pipeline (Aeon memory + CAMP/Catfish negotiation + KnowBias/RBD anti-bias), we run a controlled A/B test: 30 prompts across 6 domains, comparing the full pipeline against raw adapter-enhanced inference (no cognitive augmentation). Responses are scored on a composite metric combining keyword coverage, response length, and coherence.

\begin{table}[h]
\centering
\caption{Cognitive layer A/B comparison (composite score, higher is better). $n{=}5$ prompts per domain.}
\label{tab:cognitive-ab}
\begin{tabular}{lcccc}
\toprule
\textbf{Domain} & \textbf{Full pipeline} & \textbf{Raw (no cognitive)} & \textbf{$\Delta$} & \textbf{Winner} \\
\midrule
power & 0.554 & 0.484 & $+$0.070 & Full \\
python & 0.492 & 0.453 & $+$0.039 & Full \\
chat-fr & 0.407 & 0.397 & $+$0.010 & Full \\
spice & 0.507 & 0.509 & $-$0.002 & Tie \\
stm32 & 0.136 & 0.184 & $-$0.048 & Raw \\
embedded & 0.267 & 0.398 & $-$0.131 & Raw \\
\midrule
\textbf{Global} & 0.394 & 0.404 & $-$0.010 & 14 full / 16 raw \\
\bottomrule
\end{tabular}
\end{table}

The cognitive layer provides a clear benefit on conversational and design-oriented domains (power electronics design: $+7$\,pp, Python design patterns: $+3.9$\,pp, French conversation: $+1$\,pp) where episodic memory recall and multi-perspective negotiation add context that the base adapter alone cannot provide. However, on purely technical domains (embedded: $-13.1$\,pp, stm32: $-4.8$\,pp), the cognitive layer actively hurts performance. We hypothesize that the negotiation stage introduces unnecessary verbosity and occasionally overrides correct but terse technical answers with longer but less precise alternatives. This result motivates domain-conditional cognitive gating: the pipeline should bypass negotiation and anti-bias filtering for domains where raw adapter output is sufficient.

% ===========================================================================
% 7. EVALUATION
% ===========================================================================
\section{Evaluation}
\label{sec:eval}

\subsection{Router Accuracy}

Router V4 accuracy is reported in \Cref{tab:router}. The improvement from V3 to V4 is driven by three factors: the higher-dimensional embedding (768d vs.\ 384d), the switch from BCE to cross-entropy with class weights, and the domain merge reducing classification ambiguity from 35 to 31 classes.

Error analysis shows that remaining misclassifications concentrate on domains with limited training data (freecad: 219 examples) and domains with high vocabulary overlap (stm32 vs.\ embedded vs.\ arm).

\subsection{Adapter Health}

All 35 V4 adapters pass the health validator: no \texttt{lora\_B} weights are identically zero (a failure mode observed in pre-pivot V3 adapters). The uniform rank-16 configuration produces consistent training dynamics across domains.

\begin{table}[h]
\centering
\caption{V4 adapter training results (selected domains).}
\label{tab:adapters}
\begin{tabular}{lccc}
\toprule
\textbf{Domain} & \textbf{Examples} & \textbf{Val Loss} & \textbf{Iters} \\
\midrule
chat-fr (foundation) & 12,000 & 0.849 & 1000 \\
reasoning (foundation) & 8,500 & 0.638 & 1000 \\
python (foundation) & 6,200 & 0.712 & 1000 \\
cpp (coding) & 3,100 & 0.685 & 500 \\
kicad-dsl (niche) & 694 & 0.420 & 200 \\
spice-sim (niche) & 368 & 0.380 & 150 \\
stm32 (niche) & 711 & 0.440 & 200 \\
freecad (niche) & 219 & 0.550 & 100 \\
\bottomrule
\end{tabular}
\end{table}

\subsection{Inference Latency}

\begin{table}[h]
\centering
\caption{Pipeline latency breakdown (Mac Studio M3 Ultra, Q4).}
\label{tab:latency}
\begin{tabular}{lc}
\toprule
\textbf{Stage} & \textbf{Latency} \\
\midrule
Router (mpnet + MLP) & 8\,ms \\
LoRA swap (unpatch) & 10\,ms \\
Aeon recall (Atlas + Trace) & 62\,ms \\
MLX inference (Q4 + 8-bit KV) & 3.6--5.2\,s \\
Negotiator (CAMP + Catfish) & $<$50\,ms \\
Anti-bias (KnowBias + RBD) & $<$20\,ms \\
Aeon write & 32\,ms \\
\midrule
\textbf{Total pipeline} & \textbf{4--6\,s} \\
\bottomrule
\end{tabular}
\end{table}

The inference stage dominates total latency at 3.6--5.2\,s, which is inherent to running a 35B MoE model in Q4 quantization on Apple Silicon. All cognitive pipeline stages combined add less than 200\,ms overhead.

\subsection{Perplexity Evaluation}

We measure in-domain perplexity (PPL) on held-out validation splits for 6 representative domains, comparing the base model against the domain-adapted model.

\begin{table}[h]
\centering
\caption{In-domain perplexity (lower is better). All 6 domains show improvement.}
\label{tab:ppl}
\begin{tabular}{lcccc}
\toprule
\textbf{Domain} & \textbf{Base PPL} & \textbf{Adapted PPL} & \textbf{$\Delta$PPL} & \textbf{Improvement} \\
\midrule
chat-fr & 8.62 & 2.43 & $-$6.19 & 71.8\% \\
kicad-dsl & 6.55 & 2.54 & $-$4.01 & 61.3\% \\
kicad-pcb & 6.15 & 3.17 & $-$2.98 & 48.5\% \\
spice & 6.24 & 3.29 & $-$2.95 & 47.3\% \\
python & 4.58 & 2.05 & $-$2.53 & 55.2\% \\
reasoning & 4.30 & 2.31 & $-$1.99 & 46.2\% \\
\midrule
\textbf{Mean} & 6.07 & 2.63 & $-$3.44 & 55.1\% \\
\bottomrule
\end{tabular}
\end{table}

All adapters substantially reduce in-domain perplexity (mean $\Delta$PPL~$= -3.44$, range $-1.99$ to $-6.19$). The largest gains appear on domains furthest from the base model's pretraining distribution (chat-fr, KiCad DSL).

\subsection{Code Generation Benchmarks}

We evaluate coding adapters on three standard benchmarks: HumanEval~\citep{chen2021evaluating} ($n{=}164$, full set), MBPP ($n{=}50$), and GSM8K ($n{=}50$). HumanEval results are from a controlled sweep at temperature~0; MBPP and GSM8K are from a smaller directional sweep (standard errors in the $\pm 6$--$10$\,pp range at $n{=}50$).

\begin{table}[h]
\centering
\caption{Code generation benchmark results ($\Delta$~pass@1 vs.\ base).}
\label{tab:benchmarks}
\begin{tabular}{lcccc}
\toprule
\textbf{Adapter} & \textbf{HumanEval} ($n{=}164$) & \textbf{MBPP} ($n{=}50$) & \textbf{GSM8K} ($n{=}50$) & \textbf{Verdict} \\
\midrule
\textbf{base} & 0.604 & 0.680 & 0.520 & --- \\
cpp & $+$0.073 & $-$0.06 & $+$0.02 & Wash \\
python & $+$0.067 & $-$0.10 & $+$0.04 & Wash \\
typescript & $+$0.061 & $-$0.10 & $+$0.06 & Wash \\
rust & $+$0.043 & $-$0.08 & $+$0.18 & Anomalous \\
shell & $+$0.006 & $+$0.04 & $+$0.20 & Anomalous \\
math & $-$0.050 & $-$0.24 & $+$0.00 & \textbf{Toxic} \\
\bottomrule
\end{tabular}
\end{table}

\paragraph{HumanEval gains do not generalize to MBPP.}
Four coding adapters show consistent HumanEval uplift ($+4.3$ to $+7.3$\,pp), but all regress on MBPP ($-6$ to $-10$\,pp). The V4 adapters appear overfit to HumanEval's problem distribution --- single-function synthesis with docstring specification --- and do not transfer to MBPP's broader task format. This is a critical limitation: adapter gains are benchmark-specific rather than reflecting general coding capability improvement.

\paragraph{MBPP Python V2: stub filtering does not recover the regression.}
To test whether the MBPP regression on the Python adapter is caused by stub artifacts in the training data, we re-evaluated on a larger sample ($n{=}100$) with a cleaned V2 adapter that filters stub-only responses during dataset construction. Results: base 73\% pass@1, Python adapter (original) 66\% ($-7$\,pp), Python adapter V2 (stubs filtered) 67\% ($-6$\,pp). Stub filtering recovers only $+1$\,pp. The regression is \emph{fundamental}: it reflects the specialization-generalization tradeoff inherent to domain LoRA, not a data quality issue. The adapter compresses the domain distribution, narrowing capability on out-of-distribution MBPP tasks regardless of training set cleanliness. This confirms that the $-7$\,pp MBPP regression for the Python adapter is structural.

\paragraph{Math adapter is toxic.}
The math adapter produces $-24$\,pp on MBPP and $-5$\,pp on HumanEval while providing no uplift on GSM8K. This adapter actively destroys general coding capability and is flagged for exclusion from production serving.

\paragraph{GSM8K anomalies.}
Shell ($+20$\,pp) and Rust ($+18$\,pp) on GSM8K are suspicious given the domain mismatch. We attribute this to the small sample size ($n{=}50$): the base model achieves 0.520 on this subset vs.\ 0.455 on the full $n{=}200$ set, suggesting the subset is biased toward easier problems. These numbers require confirmation at larger $n$.

\paragraph{Cross-domain interference.}
Non-coding adapters applied to HumanEval show predictable degradation: yaml-json $-2.4$\,pp, chat-fr $-6.1$\,pp, math $-11.0$\,pp. The magnitude correlates with domain distance from Python, confirming that LoRA deltas applied uniformly across all 256 MoE experts produce inter-domain smoothing.

\subsection{System Latency}

The system latency breakdown is reported in \Cref{tab:latency}. Inference dominates at 3.6--5.2\,s; all cognitive pipeline stages combined add less than 200\,ms overhead. The 8-bit KV cache with LRU prefix reuse provides an approximate 2$\times$ speedup for sequential queries within the same domain.

\subsection{Stress Test}

The deployed system passes a full stress test: 10/10 queries routed correctly to their target domains on initial deployment, and 8/8 queries pass after a cold restart, confirming that adapter loading and routing state are properly persisted.

% ===========================================================================
% 8. DISCUSSION
% ===========================================================================
\section{Discussion}
\label{sec:discussion}

\paragraph{Contribution reframing.}
The primary contribution of \textsc{micro-kiki} is not SOTA adapter performance on any single benchmark, but rather the \emph{system-level integration} of multi-adapter routing, cognitive augmentation, and honest multi-benchmark evaluation on consumer hardware. The ablation results reveal both where the approach works (in-domain PPL, HumanEval) and where it fails (MBPP generalization, cognitive layer on technical domains). We believe that publishing negative and mixed results alongside positive ones is essential for the community to calibrate expectations around multi-LoRA systems.

\paragraph{The null-space negative result.}
The finding that null-space projection has no effect on MoE architectures (\Cref{tab:nullspace}) has practical implications: it eliminates a complex and computationally expensive training step (projecting gradients into prior adapters' null spaces) for any system built on MoE backbones with sufficient expert count. The natural quasi-orthogonality of adapter weights in high-dimensional MoE parameter spaces (${\sim}88.6$° across all tested pairs) suggests that catastrophic forgetting between adapters is structurally unlikely on these architectures, independent of training order.

\paragraph{Adapter overfitting.}
The divergence between HumanEval ($+7.3$\,pp) and MBPP ($-10$\,pp) for the same adapter is a cautionary result. Single-benchmark evaluation would have produced a misleadingly positive narrative. We attribute this to the teacher distillation process: the training data generated by Qwen3-Coder-480B-A35B may favor the structured, docstring-driven format characteristic of HumanEval problems. Multi-benchmark validation gates should be mandatory before any adapter is promoted to production.

\paragraph{Specialization-generalization tradeoff.}
A consistent pattern emerges across all evaluation axes: LoRA adapters improve \emph{in-distribution} performance substantially (PPL reduces 71\% across all 35 domains, mean $\Delta$PPL~$= -3.44$), but \emph{out-of-distribution} performance regresses for certain domains. The clearest case is the Python adapter: MBPP pass@1 drops from 73\% (base) to 66\% ($-7$\,pp), and the V2 ablation with stub filtering recovers only $+1$\,pp, confirming the regression is not a data quality artifact. This is the classical specialization-generalization tradeoff: LoRA adapters act as \emph{domain compressors}, concentrating the weight update budget on domain-specific patterns at the cost of breadth. They are not domain enhancers that uniformly improve performance on all tasks touching that domain. The practical consequence is that multi-benchmark evaluation is not merely advisable --- it is load-bearing for correctness. A system that evaluates adapters only on in-domain PPL or HumanEval would promote a Python adapter that actively degrades MBPP performance. Production mitigation is the \texttt{BASE\_ONLY\_DOMAINS} gate (\S\ref{sec:domain-gating}): when measured out-of-distribution regression exceeds 5\,pp, the adapter is bypassed and the base model serves the domain directly.

\paragraph{Cognitive layer is domain-dependent.}
The A/B evaluation (\Cref{tab:cognitive-ab}) shows that cognitive augmentation is not universally beneficial. The negotiation and anti-bias stages add value on open-ended tasks (power electronics design, French conversation) where multiple valid response strategies exist, but degrade performance on tasks with a single correct answer (register configuration, firmware implementation). This motivates a domain-conditional gating mechanism that bypasses the cognitive layer for purely technical domains.

\paragraph{Limitations.}

Router accuracy on sparse domains remains a challenge. Domains with fewer than 300 training examples (freecad, platformio) show higher misclassification rates. The current sigmoid architecture treats each domain independently; exploring cross-domain attention mechanisms may improve routing for related domains.

STM32 HAL adapter quality is limited by the teacher model's own knowledge gaps in vendor-specific APIs. Generated training data sometimes produces plausible but incorrect register configurations, which the forgetting gate does not catch (it measures interference, not factual accuracy).

The 4--6\,s end-to-end latency, while acceptable for expert consultation use cases, precludes real-time applications. This is fundamentally limited by running a 35B model on CPU/Metal rather than dedicated GPU inference.

\paragraph{Apple Silicon Specificity.}
The system's reliance on MLX and unified memory means it cannot directly port to CUDA-based deployments. However, the architectural patterns --- unpatch-based swap, sigmoid multi-label routing, cognitive pipeline staging --- are framework-agnostic.

\paragraph{Future Work.}
\begin{enumerate}[nosep]
  \item \textbf{Recursive Language Model (RLM)}: Replacing the Negotiator with a recursive decomposition model~\citep{rlm2025} for multi-domain queries. RLM decomposes a complex query into domain-specific sub-queries, routes each independently, and recursively composes the results --- a more principled approach than the current CAMP arbitration for queries spanning multiple adapter domains.
  \item \textbf{Domain-conditional cognitive gating}: Automatically bypassing negotiation and anti-bias stages for domains where the A/B evaluation shows negative $\Delta$.
  \item \textbf{Router V5}: Cross-attention routing and hard-negative mining for sparse domains.
  \item \textbf{Multi-benchmark adapter gates}: Requiring positive $\Delta$ on at least two independent benchmarks before promoting any adapter to production serving.
  \item \textbf{DPO/GRPO alignment}: Currently blocked on MLX native support for preference optimization.
  \item \textbf{Domain-specific benchmarks}: KiCad DSL correctness, SPICE netlist validity, STM32 HAL compilation rates.
  \item \textbf{Multi-node deployment}: Extending the pipeline across the existing 5-node cluster for parallel inference.
\end{enumerate}

% ===========================================================================
% 9. CONCLUSION
% ===========================================================================
\section{Conclusion}
\label{sec:conclusion}

We presented \textsc{micro-kiki}, a production system that deploys 35 domain-specialized LoRA adapters with a 7-stage cognitive pipeline on a single Mac Studio M3~Ultra. The system demonstrates that consumer Apple Silicon hardware, combined with MoE model sparsity and efficient adapter management, can serve a meaningful number of specialized domains with sub-second routing and cognitive augmentation.

Key technical contributions include: (1)~O(10\,ms) adapter swap via in-place LoRA unpatch; (2)~sigmoid multi-label routing achieving 70.7\% top-1 accuracy (2.2$\times$ over keyword baselines); (3)~the negative result that null-space projection is superfluous on native MoE architectures due to natural quasi-orthogonality at ${\sim}88.6$°; (4)~honest multi-benchmark evaluation revealing that HumanEval gains ($+7.3$\,pp) do not transfer to MBPP ($-10$\,pp); and (5)~evidence that cognitive augmentation is domain-dependent, helping design tasks but degrading technical ones. All 35 adapters reduce in-domain perplexity by 46--72\% (mean $\Delta$PPL~$= -3.4$) and exhibit no catastrophic forgetting (mean cross-stack angle: 79.4°).

The system is deployed in production and available as open-source under the Apache 2.0 license. The router is published at \texttt{clemsail/micro-kiki-router-v4} on HuggingFace.

% ===========================================================================
% REFERENCES
% ===========================================================================
\bibliographystyle{plainnat}

\begin{thebibliography}{20}

\bibitem[Hu et~al.(2022)]{hu2022lora}
Edward~J. Hu, Yelong Shen, Phillip Wallis, Zeyuan Allen-Zhu, Yuanzhi Li, Shean Wang, Lu~Wang, and Weizhu Chen.
\newblock {LoRA}: Low-Rank Adaptation of Large Language Models.
\newblock In \emph{ICLR}, 2022.
\newblock arXiv:2106.09685.

\bibitem[Gao et~al.(2024)]{gao2024mola}
Zhimeng Gao, Bosheng Ding, Chengqi Zhao, Junjie Wang, and Lidong Bing.
\newblock Higher Layers Need More {LoRA} Experts.
\newblock \emph{arXiv preprint arXiv:2403.03432}, 2024.

\bibitem[{Activated LoRA}(2024)]{activatedlora2024}
{Activated LoRA Authors}.
\newblock Activated {LoRA}: Dynamic Adapter Activation for Parameter-Efficient Fine-Tuning.
\newblock \emph{arXiv preprint arXiv:2512.17910}, 2024.

\bibitem[{LoRA Rank}(2025)]{lora_rank2025}
{LoRA Rank Authors}.
\newblock On the Optimal Rank of {LoRA}: Vanilla {LoRA} r=16 Suffices When Learning Rate Is Tuned.
\newblock \emph{arXiv preprint arXiv:2602.04998}, 2025.

\bibitem[{MobiLoRA}(2025)]{mobilora2025}
{MobiLoRA Authors}.
\newblock {MobiLoRA}: Efficient {LoRA} on Mobile Devices.
\newblock In \emph{ACL}, 2025.

\bibitem[{Aeon}(2025)]{aeon2025}
{Aeon Authors}.
\newblock Aeon: Episodic Memory with {SIMD} Vector and Neuro-Symbolic Graph Stores.
\newblock \emph{arXiv preprint arXiv:2601.15311}, 2025.

\bibitem[{KnowBias}(2025)]{knowbias2025}
{KnowBias Authors}.
\newblock {KnowBias}: Detecting and Quantifying Knowledge-Correlated Biases in {LLM} Outputs.
\newblock \emph{arXiv preprint arXiv:2601.21864}, 2025.

\bibitem[{CAMP}(2026)]{camp2026}
{CAMP Authors}.
\newblock {CAMP}: Structured Multi-Perspective Arbitration for Language Model Outputs.
\newblock \emph{arXiv preprint arXiv:2604.00085}, 2026.

\bibitem[{Catfish}(2025)]{catfish2025}
{Catfish Authors}.
\newblock Catfish: Productive Dissent Injection for Consensus-Driven {AI} Systems.
\newblock \emph{arXiv preprint arXiv:2505.21503}, 2025.

\bibitem[{RBD}(2025)]{rbd2025}
{RBD Authors}.
\newblock Residual Bias Detection as Post-Processing for Language Models.
\newblock \emph{arXiv preprint arXiv:2505.17100}, 2025.

\bibitem[Kwon et~al.(2023)]{kwon2023efficient}
Woosuk Kwon, Zhuohan Li, Siyuan Zhuang, Ying Sheng, Lianmin Zheng, Cody~Hao Yu, Joseph~E. Gonzalez, Hao Zhang, and Ion Stoica.
\newblock Efficient Memory Management for Large Language Model Serving with {PagedAttention}.
\newblock In \emph{SOSP}, 2023.

\bibitem[Sheng et~al.(2024)]{sheng2024slora}
Ying Sheng, Shiyi Cao, Dacheng Li, Coleman Hooper, Nicholas Lee, Shuo Yang, Christopher Chou, Banghua Zhu, Lianmin Zheng, Kurt Keutzer, Joseph~E. Gonzalez, and Ion Stoica.
\newblock {S-LoRA}: Serving Thousands of Concurrent {LoRA} Adapters.
\newblock In \emph{MLSys}, 2024.

\bibitem[Song et~al.(2020)]{song2020mpnet}
Kaitao Song, Xu~Tan, Tao Qin, Jianfeng Lu, and Tie-Yan Liu.
\newblock {MPNet}: Masked and Permuted Pre-training for Language Understanding.
\newblock In \emph{NeurIPS}, 2020.

\bibitem[Chen et~al.(2021)]{chen2021evaluating}
Mark Chen, Jerry Tworek, Heewoo Jun, Qiming Yuan, Henrique~Ponde de~Oliveira Pinto, Jared Kaplan, Harri Edwards, Yuri Burda, Nicholas Joseph, Greg Brockman, et~al.
\newblock Evaluating Large Language Models Trained on Code.
\newblock \emph{arXiv preprint arXiv:2107.03374}, 2021.

\bibitem[{Qwen Team}(2025)]{qwen3_2025}
{Qwen Team}.
\newblock Qwen3 Technical Report.
\newblock \emph{arXiv preprint}, 2025.

\bibitem[{Davison et~al.}(2024)]{davison2024brainstacks}
{Joe Davison, Anish Acharya, and others}.
\newblock Brainstacks: Orthogonal Adapter Stacking via Null-Space Projection.
\newblock \emph{arXiv preprint arXiv:2410.15781}, 2024.

\bibitem[{RLM Authors}(2025)]{rlm2025}
{RLM Authors}.
\newblock Recursive Language Models for Multi-Domain Decomposition.
\newblock \emph{arXiv preprint arXiv:2512.24601}, 2025.

\end{thebibliography}

\end{document}
